% ============================================================
% Chapter 10: Problem 9 — The Reciprocity Law for Algebraic Number Fields
% ============================================================

\chapter{The Reciprocity Law for Algebraic Number Fields}

\begin{partiallybox}
Hilbert's ninth problem is \textbf{solved} in its original formulation: the Artin reciprocity theorem (1920s) provides the most general reciprocity law for abelian extensions of number fields, completing the vision of class field theory. However, the far-reaching \textbf{Langlands program}---which seeks to extend reciprocity to non-abelian extensions---remains one of the deepest ongoing programs in mathematics. The problem is both historically resolved and profoundly alive.
\end{partiallybox}

\section{Historical Context}

Number theory is, at its heart, the study of the integers and their properties. Among the most fundamental questions one can ask is deceptively simple: \emph{given an integer $a$ and a prime $p$, when is $a$ a perfect square modulo $p$?} That is, for which primes $p$ does the congruence
\[
    x^2 \equiv a \pmod{p}
\]
have a solution? This question---the question of \emph{quadratic residues}---has fascinated mathematicians since antiquity, and its answer led to one of the most beautiful and surprising theorems in all of mathematics: the \textbf{law of quadratic reciprocity}.

\subsection{Quadratic Residues and the Legendre Symbol}

To make the question precise, we need a piece of notation.

\begin{definition}[Legendre symbol]
Let $p$ be an odd prime and $a$ an integer not divisible by $p$. The \textbf{Legendre symbol} is defined by
\[
    \left(\frac{a}{p}\right) = \begin{cases} \phantom{-}1 & \text{if $a$ is a quadratic residue mod $p$ (i.e., $x^2 \equiv a \pmod{p}$ has a solution),} \\ -1 & \text{if $a$ is a quadratic non-residue mod $p$.} \end{cases}
\]
\end{definition}

For example, modulo $7$, the squares are $1^2 \equiv 1$, $2^2 \equiv 4$, $3^2 \equiv 2$. So $1, 2, 4$ are quadratic residues mod $7$, and $3, 5, 6$ are non-residues. We write $\left(\frac{1}{7}\right) = \left(\frac{2}{7}\right) = \left(\frac{4}{7}\right) = 1$ and $\left(\frac{3}{7}\right) = \left(\frac{5}{7}\right) = \left(\frac{6}{7}\right) = -1$.

Computing Legendre symbols for small primes is straightforward---just check all the squares. But is there a \emph{pattern}? Is there a rule that tells us, without computation, whether $a$ is a square mod $p$?

\subsection{Euler's Conjecture}

In the 1740s, Leonhard Euler observed a striking pattern through extensive computation. He conjectured:

\begin{quote}
\emph{If $p$ and $q$ are distinct odd primes, then $\left(\frac{q}{p}\right) = 1$ whenever $p \equiv 1 \pmod{q}$, and $\left(\frac{q}{p}\right) = -1$ whenever $p$ is a primitive root mod $q$.}
\end{quote}

In modern language, Euler was saying that $\left(\frac{q}{p}\right)$ depends only on $p \pmod{q}$. This was remarkable: it said that whether $q$ is a square mod $p$ depends not on $p$ and $q$ separately, but on their relationship. But Euler could not prove it.

\subsection{Gauss's Theorem: The Law of Quadratic Reciprocity}

Carl Friedrich Gauss proved Euler's conjecture in 1783 and published his celebrated proof in his \emph{Disquisitiones Arithmeticae} (1801). He called it the \emph{theorema aureum}---the \emph{golden theorem}---and was so proud of it that he gave eight different proofs over the course of his career.

\begin{theorem}[Law of Quadratic Reciprocity]
Let $p$ and $q$ be distinct odd primes. Then
\[
    \left(\frac{p}{q}\right)\left(\frac{q}{p}\right) = (-1)^{\frac{p-1}{2} \cdot \frac{q-1}{2}}.
\]
Equivalently:
\begin{itemize}
    \item If $p \equiv 1 \pmod{4}$ or $q \equiv 1 \pmod{4}$, then $\left(\frac{p}{q}\right) = \left(\frac{q}{p}\right)$.
    \item If $p \equiv q \equiv 3 \pmod{4}$, then $\left(\frac{p}{q}\right) = -\left(\frac{q}{p}\right)$.
\end{itemize}
\end{theorem}

To this day, the quadratic reciprocity law is considered one of the most beautiful results in mathematics. It says that the question ``is $p$ a square mod $q$?'' is intimately related to the question ``is $q$ a square mod $p$?''---a symmetry that is completely unexpected. Two seemingly independent questions about two different primes turn out to be reflections of each other, linked by a simple parity condition.

Gauss's theorem also has two important supplements:
\begin{itemize}
    \item \textbf{The first supplement:} $\left(\frac{-1}{p}\right) = (-1)^{(p-1)/2}$. So $-1$ is a square mod $p$ if and only if $p \equiv 1 \pmod{4}$.
    \item \textbf{The second supplement (Gauss's lemma):} $\left(\frac{2}{p}\right) = (-1)^{(p^2-1)/8}$. So $2$ is a square mod $p$ if and only if $p \equiv \pm 1 \pmod{8}$.
\end{itemize}

Together, these results allow us to compute $\left(\frac{a}{p}\right)$ for any $a$ and $p$ using only factorization into primes.

\subsection{Higher Reciprocity Laws}

Gauss's law was not the end of the story---it was the beginning. Once mathematicians understood quadratic reciprocity, they immediately asked: \emph{is there a cubic reciprocity law? A quartic law? A law of reciprocity for higher powers?}

\paragraph{Cubic reciprocity.} The first major step was taken by Euler and Jacobi, who studied the question of whether $a \equiv x^3 \pmod{p}$ has a solution. But the full cubic reciprocity law required working not in the ordinary integers $\mathbb{Z}$, but in the ring of \emph{Eisenstein integers} $\mathbb{Z}[\omega]$, where $\omega = e^{2\pi i/3}$ is a primitive cube root of unity. This ring has the unique factorization property, and in 1844--1847, the Irish mathematician William Rowan Hamilton and the German mathematician Peter Gustav Lejeune Dirichlet proved a cubic reciprocity law in this setting.

\paragraph{Quartic (biquadratic) reciprocity.} Gauss himself worked toward a quartic reciprocity law, which involves fourth powers and requires working in the Gaussian integers $\mathbb{Z}[i]$. This was completed by Jacobi and Eisenstein in the 1840s.

\paragraph{The pattern.} A clear pattern emerged: each higher reciprocity law requires extending the number system. Quadratic reciprocity lives in $\mathbb{Z}$. Cubic reciprocity lives in $\mathbb{Z}[\omega]$. Quartic reciprocity lives in $\mathbb{Z}[i]$. The general reciprocity law for $n$th powers seems to require the $n$th cyclotomic field $\mathbb{Q}(\zeta_n)$, where $\zeta_n = e^{2\pi i/n}$.

This led to a natural and profound question: \emph{can we formulate a reciprocity law for all number fields, not just cyclotomic ones?} What is the most general reciprocity law? This was the substance of Hilbert's ninth problem.

\section{The Problem}

\begin{problem_statement}
Find the most general and natural reciprocity law for the Legendre symbols of arbitrary algebraic number fields. That is, given an algebraic number field $K$ (a finite extension of $\mathbb{Q}$), find a reciprocity law that governs the splitting of primes---extending quadratic reciprocity from $\mathbb{Z}$ to $K$.
\end{problem_statement}

Hilbert's phrasing is characteristically bold: he wanted a \emph{general} reciprocity law, one that would encompass all the known cases---quadratic, cubic, quartic---and go beyond them. The problem was not to find a formula for $\left(\frac{a}{p}\right)$ in some special case, but to discover the deep structural principle that \emph{explains why reciprocity holds at all}.

\section{Key Developments}

The resolution of Hilbert's ninth problem unfolded over the first half of the twentieth century, through the development of one of the most important theories in modern mathematics: \textbf{class field theory}.

\subsection{From Quadratic Reciprocity to Ideal Theory}

The first major conceptual advance came from Ernst Kummer (1840s--1850s) and his student Richard Dedekind (1860s--1870s). Kummer had studied ``ideal numbers'' to rescue unique factorization in cyclotomic fields (it fails in $\mathbb{Z}[\zeta_p]$ for $p \geq 23$). Dedekind refined Kummer's ideas into the modern theory of \emph{ideals}: in any ring of algebraic integers, every ideal factors uniquely into prime ideals.

This ideal-theoretic framework transformed the question. Instead of asking ``is $a$ a square mod $p$?'' (which doesn't quite make sense in a general number field), we can ask: \emph{given a prime ideal $\mathfrak{p}$ of $K$, how does it split in an extension field $L$?} Does $\mathfrak{p}$ remain prime? Does it factor into several primes? The answer depends on the relationship between $K$ and $L$, and reciprocity is the principle that governs this dependence.

\subsection{Hilbert Class Field and the Hilbert Symbol}

Hilbert himself made major contributions to this program. In the 1890s, he introduced the \emph{Hilbert class field} of a number field $K$: the maximal unramified abelian extension of $K$. This is the largest extension in which no prime of $K$ ramifies (splits ``badly'') and the Galois group is abelian. The Hilbert class field is controlled by the \emph{class group} of $K$---the group of fractional ideals modulo principal ideals---and its structure encodes deep information about $K$.

Hilbert also introduced the \emph{Hilbert symbol} $\left(\frac{a,b}{\mathfrak{p}}\right)$, a generalization of the Legendre symbol to number fields, and formulated conjectures about its properties that turned out to be precursors of the full reciprocity law.

\subsection{Takagi's Class Field Theory}

The modern theory of class field theory was largely created by Teiji Takagi in the 1920s. Takagi proved a remarkable existence theorem:

\begin{theorem}[Takagi, 1920]
For every \emph{ideal group} $H$ (containing the group of principal ideals, satisfying certain technical conditions), there exists a unique abelian extension $L/K$ such that the primes of $K$ that split completely in $L$ are precisely those lying in $H$.
\end{theorem}

This is an \emph{existence theorem}: it says that every ``reasonable'' set of splitting conditions can be achieved by some abelian extension. The correspondence between ideal groups and abelian extensions is now called the \textbf{Artin--Takagi correspondence}.

\subsection{The Artin Reciprocity Theorem}

The culmination of class field theory came with Emil Artin's reciprocity theorem in the 1920s. Artin's theorem is the precise, general reciprocity law that Hilbert had sought.

\begin{theorem}[Artin Reciprocity Theorem]
Let $K$ be a number field and $L/K$ a finite abelian extension with Galois group $G = \Gal(L/K)$. Then there exists a canonical surjective homomorphism (the \textbf{Artin map})
\[
    \left(\frac{L/K}{\cdot}\right) : I_K \to G,
\]
defined on the group $I_K$ of fractional ideals of $K$ coprime to the ramified primes, with the property that:
\begin{enumerate}
    \item A prime $\mathfrak{p}$ of $K$ splits completely in $L$ if and only if $\mathfrak{p}$ is in the kernel of the Artin map.
    \item For $\alpha \in K^\times$, the principal ideal $(\alpha)$ maps to the identity (the ``reciprocity law'': principal ideals are invisible to the Artin map).
    \item The Artin map respects the local--global principle: the local symbol at each prime agrees with the norm residue symbol.
\end{enumerate}
\end{theorem}

The reciprocity law is precisely the second property: \emph{the Artin map is trivial on principal ideals}. This is the abstract form of ``what goes around comes around''---the idele-theoretic generalization of the fact that $a$ is a square mod $p$ whenever $a$ is itself a square.

Artin's theorem is a masterpiece of abstraction. It unifies:
\begin{itemize}
    \item \textbf{Quadratic reciprocity} (the original law of Gauss).
    \item \textbf{Cubic and quartic reciprocity} (as special cases).
    \item \textbf{The splitting of primes in cyclotomic fields} (the Kronecker--Weber theorem, which says every abelian extension of $\mathbb{Q}$ is contained in a cyclotomic field).
    \item \textbf{Hilbert's symbol} and the ramification behavior of primes in any abelian extension.
\end{itemize}

\subsection{Hasse, Weil, and the Langlands Program}

Hasse contributed to the local version of class field theory (local reciprocity), while Andre Weil formulated reciprocity in terms of ideles in the 1930s, completing the picture for abelian extensions. But the question remained: \emph{what about non-abelian extensions?}

This is where the story takes a dramatic turn. In 1967, Robert Langlands wrote a famous letter to Andre Weil, outlining a vast generalization of class field theory that would encompass not just abelian, but \emph{all} extensions of number fields.

\begin{quote}
\emph{The Langlands program proposes that for every reductive group $G$ over a number field $K$, there is a correspondence between:
\begin{enumerate}
    \item \textbf{$G$-extensions} of $K$ (Galois representations with values in $G$), and
    \item \textbf{Automorphic representations} of $G$ (certain representations of the adelic group $G(\mathbb{A}_K)$).
\end{enumerate}}
\end{quote}

When $G = \mathrm{GL}_1$, the Langlands program reduces to class field theory: automorphic representations of $\mathrm{GL}_1$ are just Hecke characters, and the reciprocity law is Artin's theorem. When $G = \mathrm{GL}_2$, the program predicts deep connections between:
\begin{itemize}
    \item 2-dimensional Galois representations (Galois groups acting on vector spaces), and
    \item modular forms (analytic functions with extraordinary symmetry properties).
\end{itemize}

The proof of Fermat's Last Theorem by Andrew Wiles (1995) was essentially a proof of a special case of the Langlands program: the Taniyama--Shimura--Weil conjecture, which is the $\mathrm{GL}_2$ case. The proof of the Sato--Tate conjecture (2011) and the work on the local Langlands correspondence are further milestones. But the full Langlands program remains open, and it is now one of the grand organizing principles of modern number theory and representation theory.

\section{Current Status}

\begin{partiallybox}
The original form of Hilbert's ninth problem---find the most general reciprocity law for abelian extensions---is \textbf{solved}: the Artin reciprocity theorem provides a complete answer, and class field theory is one of the crowning achievements of 20th-century mathematics.

However, the \emph{spirit} of the problem---understand the deepest principles governing the splitting of primes in arbitrary extensions---lives on in the Langlands program. The Langlands program is very much \textbf{ongoing}, with major milestones (Fermat's Last Theorem, Sato--Tate, local Langlands) achieved but the full vision still distant. Hilbert's ninth problem is thus both historically resolved and profoundly alive.
\end{partiallybox}

\section*{Common Questions}

\begin{description}[font=\bfseries, leftmargin=2em, style=nextline]

\item[Q: What does reciprocity have to do with ``solving equations''?]\hfill\\
The most basic Diophantine equation is $x^2 = a \pmod{p}$---does $a$ have a square root modulo $p$? The Legendre symbol $\bigl(\frac{a}{p}\bigr)$ answers this question. Reciprocity laws tell us that the answer depends not on $p$ in isolation, but on $p$'s relationship to $a$ (and vice versa). Quadratic reciprocity specifically says that $\bigl(\frac{p}{q}\bigr)$ and $\bigl(\frac{q}{p}\bigr)$ are nearly the same---the question ``is $p$ a square mod $q$?'' and ``is $q$ a square mod $p$?'' are linked by a simple sign. This symmetry is not just beautiful; it gives an efficient computational method to decide solvability without enumerating squares.

\item[Q: I've heard of quadratic reciprocity---what makes it ``the most beautiful theorem''?]\hfill\\
Gauss called it the \emph{theorema aureum} (golden theorem). Its beauty lies in the unexpected symmetry it reveals: two seemingly independent questions about different primes turn out to be mirrors of each other. The statement is simple enough to explain to a beginner, yet it connects modular arithmetic, the geometry of numbers, and complex analysis. Moreover, the theorem spawns entire fields of inquiry---higher reciprocity laws, class field theory, and the Langlands program. The single equation $\bigl(\frac{p}{q}\bigr)\bigl(\frac{q}{p}\bigr) = (-1)^{\frac{p-1}{2}\frac{q-1}{2}}$ is a remarkably compact expression for a deep structural fact about prime numbers.

\item[Q: What is the Langlands program in one paragraph?]\hfill\\
The Langlands program is a vast web of conjectures proposing a fundamental correspondence between two seemingly unrelated worlds: \emph{Galois representations} (symmetries of number fields) and \emph{automorphic forms} (highly symmetric analytic functions on adele groups). When the Galois representation is one-dimensional (abelian), the correspondence reduces to class field theory---Artin reciprocity. For higher dimensions, it predicts that every $n$-dimensional Galois representation corresponds to an automorphic form on $\mathrm{GL}_n$. Fermat's Last Theorem was proved by establishing a $\mathrm{GL}_2$ case (the modularity theorem). The full program is one of the grand unifying visions of modern mathematics, but it remains largely unproven.

\item[Q: Why does cubic reciprocity need Eisenstein integers instead of ordinary integers?]\hfill\\
Cubic reciprocity asks when $a$ is a cube modulo $p$, i.e., does $x^3 \equiv a \pmod{p}$ have a solution? To find a symmetry analogous to quadratic reciprocity, you need to work in a number system where cube roots of unity exist---the Eisenstein integers $\mathbb{Z}[\omega]$ where $\omega = e^{2\pi i/3}$. The reason is that the cubic character $\chi(a)$ naturally takes values that are powers of $\omega$, not $\pm 1$, and the reciprocity law involves a pairing that only makes sense in the larger ring. This pattern repeats for every higher power: $n$th power reciprocity requires adjoining $\zeta_n$, an $n$th root of unity.

\item[Q: What is the Artin map in simple terms?]\hfill\\
The Artin map is a rule that assigns to each prime ideal $\mathfrak{p}$ of a number field $K$ an element of the Galois group of an abelian extension $L/K$. This element---the \emph{Frobenius} element---encodes how the prime $\mathfrak{p}$ behaves in $L$: whether it splits, remains inert, or ramifies. The deep part of Artin reciprocity is that this map factors through the ideal class group: principal ideals (generated by a single element) all map to the identity. This means that the splitting behavior of a prime depends only on its ideal class, not on the prime itself---a sweeping generalization of quadratic reciprocity.

\item[Q: What does ``abelian extension'' mean and why is it important for reciprocity?]\hfill\\
An extension $L/K$ is \emph{abelian} if its Galois group (the group of field automorphisms fixing $K$) is abelian, meaning its elements commute. This condition is strong enough to force the extension to have a very structured form. Class field theory gives a complete description of all abelian extensions of any number field---they correspond to certain groups of ideals. Reciprocity laws like quadratic and cubic reciprocity are special cases of this correspondence. Non-abelian extensions (like the splitting field of $x^3 - 2$) are far more complex and are the subject of the Langlands program.

\item[Q: Is the Langlands program proven or still a conjecture?]\hfill\\
The Langlands program is a vast collection of conjectures, some of which have been proved and many of which remain open. The $\mathrm{GL}_1$ case is class field theory (proved). The $\mathrm{GL}_2$ case over $\mathbb{Q}$ was largely proved by the modularity theorem (Wiles et al., 1995) and subsequent work. The local Langlands correspondence for $\mathrm{GL}_n$ over $p$-adic fields was proved by Harris--Taylor and Henniart (2000). However, the full global Langlands correspondence for arbitrary reductive groups over arbitrary number fields remains one of the deepest open problems in mathematics. Progress continues steadily, but a complete proof is likely decades away.

\end{description}

\section*{Exercises}

\begin{enumerate}
    \item \textbf{Computing Legendre symbols.}
    \begin{enumerate}
        \item Compute $\left(\frac{2}{7}\right)$, $\left(\frac{3}{7}\right)$, and $\left(\frac{5}{7}\right)$ by listing all squares mod~$7$.
        \item Compute $\left(\frac{5}{11}\right)$ and $\left(\frac{11}{5}\right)$ directly, and verify the quadratic reciprocity law for the pair $(5, 11)$.
        \item Compute $\left(\frac{7}{11}\right)$ and $\left(\frac{11}{7}\right)$ directly, and verify quadratic reciprocity for the pair $(7, 11)$.
    \end{enumerate}

    \item \textbf{Using quadratic reciprocity.} Use the law of quadratic reciprocity and the two supplements to determine whether each of the following congruences has a solution:
    \begin{enumerate}
        \item $x^2 \equiv 3 \pmod{13}$
        \item $x^2 \equiv 5 \pmod{29}$
        \item $x^2 \equiv 7 \pmod{127}$
        \item $x^2 \equiv -1 \pmod{53}$
    \end{enumerate}

    \item \textbf{Euler's criterion.} \emph{Euler's criterion} states that for an odd prime $p$ and $\gcd(a,p) = 1$,
    \[
        \left(\frac{a}{p}\right) \equiv a^{(p-1)/2} \pmod{p}.
    \]
    \begin{enumerate}
        \item Use Euler's criterion to compute $\left(\frac{3}{13}\right)$.
        \item Prove Euler's criterion (hint: consider the map $x \mapsto x^2$ on $(\mathbb{Z}/p\mathbb{Z})^\times$ and use Fermat's little theorem).
    \end{enumerate}

    \item \textbf{The Jacobi symbol.} The \textbf{Jacobi symbol} extends the Legendre symbol to composite moduli. If $n = p_1^{a_1} \cdots p_k^{a_k}$ is odd and positive, define
    \[
        \left(\frac{a}{n}\right) = \prod_{i=1}^{k} \left(\frac{a}{p_i}\right)^{a_i}.
    \]
    \begin{enumerate}
        \item Compute $\left(\frac{2}{35}\right)$ and $\left(\frac{3}{35}\right)$.
        \item Show that the Jacobi symbol satisfies the same multiplicative properties as the Legendre symbol: $\left(\frac{ab}{n}\right) = \left(\frac{a}{n}\right)\left(\frac{b}{n}\right)$ and $\left(\frac{a}{mn}\right) = \left(\frac{a}{m}\right)\left(\frac{a}{n}\right)$.
        \item \textbf{Warning:} Does $\left(\frac{a}{n}\right) = 1$ imply that $a$ is a square mod $n$? Check: is $2$ a square mod $35$? (Hint: compute $\left(\frac{2}{5}\right)$ and $\left(\frac{2}{7}\right)$ separately.)
    \end{enumerate}

    \item \textbf{Gauss's Lemma.} \emph{Gauss's lemma} states: let $p$ be an odd prime and $a$ an integer with $\gcd(a,p) = 1$. Consider the residues $a, 2a, 3a, \ldots, \frac{p-1}{2}a$ reduced modulo $p$ to the range $\{-\frac{p-1}{2}, \ldots, -1, 1, \ldots, \frac{p-1}{2}\}$. Let $n$ be the number of these residues that are negative. Then $\left(\frac{a}{p}\right) = (-1)^n$.
    \begin{enumerate}
        \item Use Gauss's lemma to compute $\left(\frac{2}{7}\right)$.
        \item Use Gauss's lemma to compute $\left(\frac{3}{7}\right)$.
        \item Prove Gauss's lemma (hint: pair each residue $r$ with $-r$ in the product $\prod_{k=1}^{(p-1)/2} (ka)$).
    \end{enumerate}

    \item $\star$ \textbf{What does ``generalization'' mean?} Hilbert asked for a reciprocity law for ``arbitrary number fields.'' The Artin reciprocity theorem provides such a law for abelian extensions. But the Langlands program goes further, seeking reciprocity for non-abelian extensions.
    \begin{enumerate}
        \item Explain what it means for a Galois extension to be ``abelian'' versus ``non-abelian.'' Give one example of each.
        \item In quadratic reciprocity, the Legendre symbol $\left(\frac{a}{p}\right)$ takes values in $\{+1, -1\}$, which is the Galois group of $\mathbb{Q}(\sqrt{a})/\mathbb{Q}$ when it is abelian. How does this relate to the Artin map?
        \item Why is it ``harder'' to study non-abelian extensions? (Hint: think about what information the Artin map carries when the Galois group is abelian versus non-abelian.)
    \end{enumerate}

    \item $\star$ \textbf{The Kronecker--Weber theorem.} The Kronecker--Weber theorem states that every finite abelian extension of $\mathbb{Q}$ is contained in a cyclotomic field $\mathbb{Q}(\zeta_n)$ for some $n$.
    \begin{enumerate}
        \item The field $\mathbb{Q}(\sqrt{5})$ is an abelian extension of $\mathbb{Q}$. Find an $n$ such that $\mathbb{Q}(\sqrt{5}) \subseteq \mathbb{Q}(\zeta_n)$. (Hint: what is the conductor?)
        \item The field $\mathbb{Q}(\sqrt[3]{2})$ is \emph{not} abelian over $\mathbb{Q}$. Why does the Kronecker--Weber theorem not apply?
    \end{enumerate}
\end{enumerate}
